\chapter{第17套试卷——参考答案与评分标准}

\section{一、选择题答案（18分）}

\begin{enumerate}[label=\arabic*.]
\item \textbf{B} — CPLD基于乘积项；FPGA基于LUT。A混淆结构，C错(FPGA SRAM掉电丢失)，D错(CPLD非易失)。
\item \textbf{C} — CASE并行无优先级；IF-ELSE级联有优先级。A错(IF-ELSE不一定优先)，B错(误导)，D与实际相反。
\item \textbf{B} — COMPONENT声明在ARCHITECTURE声明区(IS与BEGIN之间)。
\item \textbf{D} — GENERIC语法：\texttt{name : type := default}，用\texttt{:=}不是\texttt{<=}。
\item \textbf{A} — FOR GENERATE循环复制；IF GENERATE条件生成。
\item \textbf{D} — One-Hot(N个FF译码简单)；Binary($\lceil\log_2 N\rceil$个FF译码复杂)。
\end{enumerate}

\section{二、填空题答案（12分）}

\begin{enumerate}[label=\arabic*.]
\item \textbf{（1）COMPONENT}在声明区，\textbf{（2）PORT MAP}实例化连接
\item \textbf{（1）GENERIC}在PORT\textbf{之前}，\textbf{（2）GENERIC MAP}传参
\item \textbf{（1）FOR GENERATE}（重复），\textbf{（2）IF GENERATE}（条件）
\item One-Hot: N个FF；Binary: $\lceil\log_2 N\rceil$个FF
\end{enumerate}

\section{三、程序改错题解析（5分）}
{\small
\begin{enumerate}[label=错误\arabic*：]
\item COMPONENT名称与实例化标签混用
\item PORT MAP进位索引错误(\texttt{carry(i-1)}应为\texttt{carry(i)})
\item 重复的ENTITY声明块
\item 敏感列表不完整
\item 进位链索引越界(carry位宽不足)
\end{enumerate}
}

\section{四、程序阅读分析题解析（5分）}

\textbf{（1）（1分）}模60 BCD计数器（个位0-9，十位0-5）。\\
\textbf{（2）（1分）}异步高电平复位(\texttt{if rst='1'})。\\
\textbf{（3）（1.5分）}$tc=1$当十位=5且个位=9（59）。\\
\textbf{（4）（1.5分）}改清零条件为$tens=9, units=9$即模100。

\section{五、综合设计题参考答案（60分）}

\subsection{设计题1：4位比较器（20分）}

\textbf{【VHDL】}
\begin{lstlisting}
library IEEE;
use IEEE.STD_LOGIC_1164.all;

entity comparator_4bit is
  port(A, B : in std_logic_vector(3 downto 0);
       AgtB, AeqB, AltB : out std_logic);
end comparator_4bit;

architecture behav of comparator_4bit is
begin
  process(A, B)
  begin
    if A > B then AgtB<='1'; AeqB<='0'; AltB<='0';
    elsif A = B then AgtB<='0'; AeqB<='1'; AltB<='0';
    else AgtB<='0'; AeqB<='0'; AltB<='1'; end if;
  end process;
end behav;
\end{lstlisting}

\textbf{【评分】（20分）}ENTITY 3分 + 三种互斥输出5分 + IF完整4分 + 敏感列表2分 + 规范2分 + 输出互斥保证4分。

\subsection{设计题2：模60 BCD计数器（20分）}

\textbf{【VHDL】}
\begin{lstlisting}
library IEEE;
use IEEE.STD_LOGIC_1164.all;
use IEEE.NUMERIC_STD.all;

entity bcd_mod60 is
  port(clk, rst : in std_logic;
       units, tens : out std_logic_vector(3 downto 0);
       tc : out std_logic);
end bcd_mod60;

architecture behav of bcd_mod60 is
  signal u_cnt, t_cnt : unsigned(3 downto 0);
begin
  process(clk, rst)
  begin
    if rst='1' then u_cnt<=(others=>'0'); t_cnt<=(others=>'0');
    elsif rising_edge(clk) then
      if t_cnt=5 and u_cnt=9 then
        u_cnt<=(others=>'0'); t_cnt<=(others=>'0');
      elsif u_cnt=9 then u_cnt<=(others=>'0'); t_cnt<=t_cnt+1;
      else u_cnt<=u_cnt+1; end if;
    end if;
  end process;
  units<=std_logic_vector(u_cnt); tens<=std_logic_vector(t_cnt);
  tc<='1' when t_cnt=5 and u_cnt=9 else '0';
end behav;
\end{lstlisting}

\textbf{【评分】（20分）}BCD进位8分 + 清零(59→00)4分 + 复位3分 + tc 2分 + NUMERIC\_STD 2分 + 端口1分。

\subsection{设计题3：密码锁控制器（20分）}

\textbf{【分析】}16位密码"2-5-3-7"=0010-0101-0011-0111，din逐位输入。18状态Moore型（S0=IDLE, S1-S16逐位匹配, S17=UNLOCK）。任意位错误→回S0。

\textbf{【VHDL】}
\begin{lstlisting}
library IEEE;
use IEEE.STD_LOGIC_1164.all;

entity combo_lock is
  port(clk, rst, din : in std_logic; unlock : out std_logic);
end combo_lock;

architecture moore_fsm of combo_lock is
  type state_type is (S0,S1,S2,S3,S4,S5,S6,S7,S8,
                      S9,S10,S11,S12,S13,S14,S15,S16,S17);
  signal curr, nxt : state_type;
begin
  reg_proc: process(clk, rst)
  begin
    if rst='1' then curr<=S0;
    elsif rising_edge(clk) then curr<=nxt; end if;
  end process;

  comb_proc: process(curr, din)
  begin
    case curr is
      when S0=> if din='0' then nxt<=S1; else nxt<=S0; end if;
      when S1=> if din='0' then nxt<=S2; else nxt<=S0; end if;
      when S2=> if din='1' then nxt<=S3; else nxt<=S0; end if;
      when S3=> if din='0' then nxt<=S4; else nxt<=S0; end if;
      when S4=> if din='0' then nxt<=S5; else nxt<=S0; end if;
      when S5=> if din='1' then nxt<=S6; else nxt<=S0; end if;
      when S6=> if din='0' then nxt<=S7; else nxt<=S0; end if;
      when S7=> if din='1' then nxt<=S8; else nxt<=S0; end if;
      when S8=> if din='0' then nxt<=S9; else nxt<=S0; end if;
      when S9=> if din='0' then nxt<=S10; else nxt<=S0; end if;
      when S10=> if din='1' then nxt<=S11; else nxt<=S0; end if;
      when S11=> if din='1' then nxt<=S12; else nxt<=S0; end if;
      when S12=> if din='0' then nxt<=S13; else nxt<=S0; end if;
      when S13=> if din='1' then nxt<=S14; else nxt<=S0; end if;
      when S14=> if din='1' then nxt<=S15; else nxt<=S0; end if;
      when S15=> if din='1' then nxt<=S16; else nxt<=S0; end if;
      when S16=> if din='1' then nxt<=S17; else nxt<=S0; end if;
      when S17=> nxt<=S17;
      when others=> nxt<=S0;
    end case;
  end process;
  unlock<='1' when curr=S17 else '0';
end moore_fsm;
\end{lstlisting}

\textbf{【评分】（20分）}状态定义4分 + 逐位匹配逻辑8分 + 错误回S0 3分 + Moore输出2分 + 代码完整3分。

\subsection{设计题3补充：状态转移图}
\begin{center}
\begin{tikzpicture}[->,>=stealth,node distance=2cm,auto,scale=0.7]
\node[draw,circle] (S0) at (0,0) {S0};
\node[draw,circle] (S1) at (1.5,0) {S1};
\node[draw,circle] (S2) at (3,0) {S2};
\node[draw,circle] (S3) at (4.5,0) {S3};
\node[draw,circle] (S4) at (6,0) {S4};
\node[draw,circle] (S17) at (3,-2) {S17/1};
\draw (S0) edge node{0} (S1);
\draw (S1) edge node{0} (S2);
\draw (S2) edge node{1} (S3);
\draw (S3) edge node{0} (S4);
\node at (7.5,0) {$\cdots$};
\draw[->] (S4) -- (7.5,0);
\draw[->] (9.5,0) -- (S17);
\draw (S17) edge[loop below] node{$\times$} (S17);
\node at (4.5,-3) {\small 每步din错误→回S0（省略未画出）};
\end{tikzpicture}
\end{center}

\textbf{【评分】（20分）}状态定义4分 + 逐位匹配8分 + 错误回S0 3分 + Moore输出2分 + 代码/图3分。
